\chapter{Tensor metryczny, długość i miara}
\label{ch:metryka-riemannowska}

Topologia mówi, które punkty są blisko w sensie otoczeń,
ale sama nie podaje długości odcinka ani pola powierzchni.
Na gładkiej rozmaitości możemy w każdym punkcie mierzyć
wektory styczne iloczynem skalarnym. Jeżeli taki pomiar
zmienia się gładko od punktu do punktu, otrzymujemy
\emph{tensor metryczny}. Z niego wyprowadzimy długość
krzywej, odległość punktów, pole podrozmaitości i objętość.
Na końcu połączymy te pojęcia z miarą Hausdorffa
i twierdzeniem Gaussa o dywergencji.

\section{Metryka w przestrzeniach stycznych}
\label{sec:metryka-definicja}

\begin{definicja}[Metryka riemannowska]
\index{metryka riemannowska@Metryka riemannowska}
\label{def:metryka-riemannowska}
Na gładkiej rozmaitości $M^n$ \emph{metryka riemannowska}
$g$ jest gładkim polem symetrycznych tensorów typu $(0,2)$
takim, że dla każdego $p\in M$ i $v\in T_pM\setminus\{0\}$
zachodzi $g_p(v,v)>0$. Piszemy
\[
 \langle v,w\rangle_{g,p}:=g_p(v,w),
 \qquad |v|_g:=\sqrt{g_p(v,v)}.
\]
Dla niezerowych $v,w$ kąt $\alpha$ wyznacza
$\cos\alpha=g_p(v,w)/(|v|_g|w|_g)$.
\end{definicja}

Intuicja: $g_p$ jest linijką i kątomierzem przyłożonym
do \emph{wektorów stycznych w punkcie $p$}. Nie porównuje
jeszcze bezpośrednio dwóch odległych punktów. Żeby to zrobić,
zsumujemy później długości maleńkich wektorów prędkości
wzdłuż krzywej. Współrzędne mogą deformować rysunek:
okrąg jednostkowych wektorów $g_p(v,v)=1$ bywa na kartce
elipsą. To zmiana opisu, nie zmiana pomiaru geometrycznego.

W mapie $x=(x^1,\ldots,x^n)$ zapisujemy
\begin{equation}\label{eq:metryka-macierz}
 g=\sum_{i,j=1}^n g_{ij}(x)\,\dd x^i\otimes\dd x^j,
 \qquad g_{ij}=g\!\left(\frac{\partial}{\partial x^i},
                          \frac{\partial}{\partial x^j}\right).
\end{equation}
Macierz $G=(g_{ij})$ jest symetryczna i dodatnio określona.
Zapis $g=g_{ij}\dd x^i\dd x^j$ jest skrótem
dla \eqref{eq:metryka-macierz}: nie chodzi tu o iloczyn
zewnętrzny form.

\begin{stwierdzenie}[Zmiana współrzędnych]
\label{prop:metryka-zmiana}
Jeżeli $x=x(y)$ na przecięciu dwóch map i
$J=(\partial x^i/\partial y^a)$, to
\[
 G_y=J^{\mathsf T}G_xJ,\qquad
 \det G_y=(\det J)^2\det G_x.
\]
\end{stwierdzenie}
\begin{proof}
Reguła łańcuchowa daje
$\partial/\partial y^a
 =\sum_i(\partial x^i/\partial y^a)\partial/\partial x^i$.
Wstawmy tę równość dwa razy do
$g(\partial/\partial y^a,\partial/\partial y^b)$.
Otrzymamy
$(G_y)_{ab}=\sum_{i,j}J_{ia}(G_x)_{ij}J_{jb}$,
czyli pierwszy wzór. Wyznacznik iloczynu macierzy
i $\det J^{\mathsf T}=\det J$ dają drugi.
W szczególności dodatnia określoność nie zależy
od wyboru współrzędnych: $w^{\mathsf T}G_yw
=(Jw)^{\mathsf T}G_x(Jw)>0$ dla $w\ne0$.
\end{proof}

\subsection*{Cofnięcie i przeniesienie metryki}
\label{subsec:metryka-pullback-pushforward}

Zapis $f^*h$ jest szczególnym przypadkiem cofnięcia
tensora z Definicji~\ref{def:cofniecie-tensora}.
Warto tu rozdzielić dwie czynności: \emph{cofnięcie}
metryki z przestrzeni docelowej na dziedzinę oraz
\emph{przeniesienie} metryki z dziedziny na cel przez
dyfeomorfizm. Nazwy angielskie to odpowiednio
\emph{pullback} i \emph{pushforward}.

\begin{definicja}[Pullback i pushforward metryki]
\index{pullback i pushforward metryki@Pullback i pushforward metryki}
\label{def:metryka-pullback-pushforward}
Niech $f:M\to N$ będzie gładkim odwzorowaniem, a $h$
metryką riemannowską na $N$. \emph{Cofnięty tensor}
$f^*h$ na $M$ jest określony przez
\begin{equation}\label{eq:metryka-pullback}
 (f^*h)_p(v,w)
 =h_{f(p)}(\dd f_pv,\dd f_pw),
 \qquad v,w\in T_pM.
\end{equation}
Jeśli $f$ jest immersją, $f^*h$ jest metryką
riemannowską, zwaną metryką indukowaną.

Jeżeli $f$ jest dyfeomorfizmem i $g$ jest metryką
na $M$, definiujemy \emph{przeniesioną metrykę}
na $N$ wzorem
\begin{equation}\label{eq:metryka-pushforward}
 (f_*g)_q(u,w)
 :=g_{f^{-1}(q)}(\dd(f^{-1})_q u,
                    \dd(f^{-1})_q w),
 \qquad f_*g=(f^{-1})^*g.
\end{equation}
\end{definicja}

W~\eqref{eq:metryka-pullback} dwa wektory z $T_pM$
przenosimy różniczką do $T_{f(p)}N$ i mierzymy tam
za pomocą $h$. Dlatego do cofnięcia tensora typu
$(0,2)$ \emph{nie} jest potrzebna odwrotność $f$.
Jeżeli $\dd f_p$ ma jądro, wektor $v\ne0$ z tego
jądra ma $(f^*h)_p(v,v)=0$; wtedy cofnięty tensor
nie jest metryką riemannowską. Na przykład dla
$f:\R^2\to\R$, $f(x,y)=x$, cofnięcie $\dd t^2$
to $\dd x^2$, które nie mierzy wektora $\partial_y$.
Przy pushforwardzie trzeba zaś umieć odtworzyć
jednoznacznie wektor w punkcie $f^{-1}(q)$;
dyfeomorfizm to gwarantuje.

\begin{stwierdzenie}[Co daje dyfeomorfizm?]
\label{prop:metryka-dyfeo-izometria}
Niech $f:M\to N$ będzie dyfeomorfizmem.
Wtedy $f^*h$ i $f_*g$ są metrykami riemannowskimi,
a operacje są wzajemnie odwrotne:
\[
 f^*(f_*g)=g,\qquad f_*(f^*h)=h.
\]
Dyfeomorfizm $f:(M,g)\to(N,h)$ jest izometrią
wtedy i tylko wtedy, gdy $f^*h=g$, równoważnie
$f_*g=h$. W tym przypadku zachowuje długości krzywych,
odległość punktów, kąty, geodezyjne i tensor krzywizny.
\end{stwierdzenie}
\begin{proof}
Jeśli $v\ne0$, to $\dd f_pv\ne0$, bo różniczka
dyfeomorfizmu jest odwracalna. Zatem
$(f^*h)_p(v,v)=h_{f(p)}(\dd f_pv,\dd f_pv)>0$;
gładkość i symetria wynikają ze wzoru.
Analogicznie jest dla $f_*g$. Równość
$\dd(f^{-1})_{f(p)}\circ\dd f_p=\mathrm{Id}$ daje
\[
 (f^*f_*g)_p(v,w)
 =(f_*g)_{f(p)}(\dd f_pv,\dd f_pw)
 =g_p(v,w).
\]
Po zamianie $f$ z $f^{-1}$ otrzymujemy drugą
tożsamość. Warunek $f^*h=g$ mówi dokładnie, że
$\dd f_p$ zachowuje iloczyn skalarny wszystkich
wektorów stycznych; jest to definicja izometrii.
Dla każdej odcinkami gładkiej krzywej $\gamma$
reguła łańcuchowa daje
\[
 L_h(f\circ\gamma)
 =\int\sqrt{h_{f(\gamma(t))}
             (\dd f\,\dot\gamma,\dd f\,\dot\gamma)}\,\dd t
 =\int\sqrt{g_{\gamma(t)}
             (\dot\gamma,\dot\gamma)}\,\dd t
 =L_g(\gamma).
\]
Ponieważ $f$ i $f^{-1}$ przenoszą wszystkie drogi
między odpowiednimi punktami, branie kresu dolnego
ich długości daje $d_h(f(p),f(q))=d_g(p,q)$.
Twierdzenia o zachowaniu geodezyjnych i krzywizny są
zapowiedzią pojęć z dwóch następnych rozdziałów.
Poniższy akapit można przeczytać po poznaniu koneksji
Leviego-Civity i definicji krzywizny; pozostała część
obecnego rozdziału nie korzysta z tych dwóch własności.
Sprawdźmy je, używając wskazanych tam definicji. Dla pól na $M$ połóżmy
\[
 \widetilde\nabla_XY
 :=f^{-1}_*\bigl(\nabla^h_{f_*X}(f_*Y)\bigr).
\]
Ponieważ $f_*[X,Y]=[f_*X,f_*Y]$, torsja tej
koneksji jest zerowa. Zgodność $\nabla^h$ z $h$
oraz $g(Y,Z)=h(f_*Y,f_*Z)\circ f$ daje
\[
 Xg(Y,Z)=g(\widetilde\nabla_XY,Z)
             +g(Y,\widetilde\nabla_XZ).
\]
Z jednoznaczności w Twierdzeniu~\ref{thm:levi-civita}
wynika $\widetilde\nabla=\nabla^g$. Stąd
$f_*(\nabla^g_{\dot\gamma}\dot\gamma)
=\nabla^h_{(f\circ\gamma)'}(f\circ\gamma)'$,
więc obrazy geodezyjnych są geodezyjnymi.
Podstawienie tej równości do definicji~\eqref{eq:riemann-def}
pokazuje, że $\dd f_p$ przenosi też operator krzywizny.
\end{proof}

W mapach, jeśli $J_f$ jest macierzą różniczki $f$,
to~\eqref{eq:metryka-pullback} przyjmuje postać
\[
 G_{f^*h}(x)=J_f(x)^{\mathsf T}H(f(x))J_f(x).
\]
W szczególności wzór
$G_y=J^{\mathsf T}G_xJ$ ze
Stwierdzenia~\ref{prop:metryka-zmiana} jest
cofnięciem \emph{tej samej} metryki przy zmianie
współrzędnych. Współczynniki się zmieniają, lecz
odległości mierzone na rozmaitości pozostają te same.

\begin{przyklad}[Dwa kierunki przekształcenia]
\label{prz:metryka-pull-push-przyklad}
Dyfeomorfizm $f:\R^2_{u,v}\to\R^2_{x,y}$,
$f(u,v)=(u,v+u^2)$, ma odwrotność
$f^{-1}(x,y)=(x,y-x^2)$. Cofając metrykę
euklidesową przestrzeni docelowej, znajdujemy
\[
 f^*(\dd x^2+\dd y^2)
 =\dd u^2+(\dd v+2u\,\dd u)^2.
\]
Przenosząc w przeciwną stronę metrykę euklidesową
dziedziny, dostajemy
\[
 f_*(\dd u^2+\dd v^2)
 =\dd x^2+(\dd y-2x\,\dd x)^2.
\]
To \emph{dwie różne} metryki zapisane na dwóch
różnych dziedzinach. Pierwsza sprawia, że
$f:(\R^2,f^*h_{\mathrm E})\to(\R^2,h_{\mathrm E})$
jest izometrią; druga sprawia, że
$f:(\R^2,g_{\mathrm E})\to(\R^2,f_*g_{\mathrm E})$
jest izometrią. Choć współczynniki zależą od punktu,
obie są płaskie, bo powstały przez przeniesienie
metryki euklidesowej.
\end{przyklad}

\begin{twierdzenie}[Istnienie metryki riemannowskiej]
\label{thm:istnienie-metryki}
Każda gładka rozmaitość Hausdorffa o przeliczalnej bazie,
również z brzegiem, ma metrykę riemannowską.
\end{twierdzenie}
\begin{proof}
Wybierzmy pokrycie mapami $(U_\alpha,x_\alpha)$.
W każdej mapie przenieśmy standardowy iloczyn skalarny
z $\R^n$ na wektory styczne:
\[
 (g_\alpha)_p(v,w)
 =\sum_{i=1}^n
   \dd x_\alpha^i{}_p(v)\,\dd x_\alpha^i{}_p(w).
\]
To dodatnio określony, gładki tensor na $U_\alpha$,
bo $\dd x_{\alpha,p}:T_pM\to\R^n$ jest izomorfizmem.
Twierdzenie~\ref{thm:podzial-jednosci-formy}
daje gładki podział jedności $(\rho_j)$
z $\supp\rho_j\subset U_{\alpha(j)}$.
Każdy tensor $\rho_jg_{\alpha(j)}$ przedłużamy przez
zero poza odpowiadającą mu mapę i kładziemy
$g=\sum_j\rho_jg_{\alpha(j)}$.
Suma jest lokalnie skończona, zatem gładka.
Jeżeli $v\ne0$ w punkcie $p$, to dla wszystkich $j$
składniki $\rho_j(p)(g_{\alpha(j)})_p(v,v)$ są
nieujemne, a przynajmniej jeden jest dodatni:
$\sum_j\rho_j(p)=1$, więc istnieje $j$ z
$\rho_j(p)>0$. Stąd $g_p(v,v)>0$.
Ta sama konstrukcja działa w mapach przybrzegowych.
\end{proof}

Twierdzenie zapewnia \emph{istnienie}, nie wskazuje
jednego wyróżnionego wyboru. Na przykład na $\R^n$, $n\geq1$,
metryki $\sum_i(\dd x^i)^2$ i
$e^{2f(x)}\sum_i(\dd x^i)^2$ są różne dla
niezerowej funkcji gładkiej $f$. Druga mnoży
długości wektorów w $x$ przez $e^{f(x)}$,
a objętości $n$-wymiarowe przez $e^{nf(x)}$.

\section{Geometrie modelowe}
\label{sec:metryki-modelowe}

\begin{przyklad}[Przestrzeń euklidesowa]
Na $\R^n$ bierzemy $g_{\mathrm E}=\sum_i(\dd x^i)^2$.
Wtedy $G=I_n$ i $|v|_g=\sqrt{\sum_i(v^i)^2}$.
To punkt odniesienia dla kolejnych modeli.
\end{przyklad}

\begin{przyklad}[Sfera]
\label{prz:metryka-sfery}
Dla $R>0$ sfera $S_R^2=\{x\in\R^3:|x|=R\}$ otrzymuje
metrykę przez ograniczenie euklidesowego iloczynu
skalarnego do $T_pS_R^2$. Parametryzacja
\[
 F(\theta,\phi)
 =R(\sin\theta\cos\phi,\sin\theta\sin\phi,\cos\theta)
\]
ma wektory pochodne
\[
 F_\theta
 =R(\cos\theta\cos\phi,\cos\theta\sin\phi,-\sin\theta),
 \quad
 F_\phi
 =R(-\sin\theta\sin\phi,\sin\theta\cos\phi,0).
\]
Ich iloczyny skalarne wynoszą kolejno
$R^2$, $0$ i $R^2\sin^2\theta$. Zatem
\begin{equation}\label{eq:metryka-sfery}
 g_{S_R^2}=R^2\dd\theta^2
       +R^2\sin^2\theta\,\dd\phi^2.
\end{equation}
Długość przy zmianie samego kąta biegunowego
$\theta$ jest inna niż przy zmianie $\phi$ blisko
bieguna: równoleżniki stają się tam krótkie.
Zanik $\sin\theta$ na biegunie nie oznacza
degeneracji metryki; współrzędna $\phi$ nie jest
tam poprawną współrzędną lokalną.

Dla $S_R^n$ współrzędne stereograficzne $u\in\R^n$
z mapą odwrotną
\[
 s(u)=\left(\frac{2R^2u}{R^2+|u|^2},
            R\frac{|u|^2-R^2}{R^2+|u|^2}\right)
\]
dają
\begin{equation}\label{eq:metryka-stereograficzna}
 s^*g_{S_R^n}
 =\left(\frac{2R^2}{R^2+|u|^2}\right)^2
   \sum_i(\dd u^i)^2.
\end{equation}
Rzeczywiście, dla $q=R^2+|u|^2$ i wektora $h$
różniczka dwóch składowych $s$ ma postać
\[
 \dd s_u(h)=
 \left(\frac{2R^2h}{q}
       -\frac{4R^2u(u\cdot h)}{q^2},
       \frac{4R^3(u\cdot h)}{q^2}\right).
\]
W kwadracie normy wyrazy zawierające
$(u\cdot h)^2$ znoszą się, a zostaje
$4R^4|h|^2/q^2$. Spolaryzowanie tej
równości daje \eqref{eq:metryka-stereograficzna}.
\end{przyklad}

\begin{przyklad}[Trzy modele przestrzeni hiperbolicznej]
\label{prz:trzy-modele-hiperboliczne}
Niech $n\geq2$. W \emph{modelu półprzestrzeni}
punkty mają współrzędne $(x,y)\in\R^{n-1}\times(0,\infty)$,
a metryka jest
\begin{equation}\label{eq:hip-polprzestrzen}
 g_{\mathbb H}=\frac{|\,\dd x\,|^2+\dd y^2}{y^2}.
\end{equation}
Przy małej wysokości $y$ ten sam ruch euklidesowy
staje się długi hiperbolicznie.

W \emph{modelu kuli Poincarégo} $B^n=\{|u|<1\}$,
\begin{equation}\label{eq:hip-kula}
 g_B=\frac{4\sum_i(\dd u^i)^2}{(1-|u|^2)^2}.
\end{equation}
Brzeg kuli nie należy do przestrzeni. Wreszcie
\emph{model hiperboloidy} jest górnym płatem
\[
 \mathcal H^n=\{(t,z)\in\R\times\R^n:
                    t^2-|z|^2=1,\ t>0\}
\]
z ograniczeniem formy $-\dd t^2+\sum_i(\dd z^i)^2$
do wektorów stycznych. Choć forma w otoczeniu
ma znak minus, jej ograniczenie do $T\mathcal H^n$
jest dodatnio określone. Jeśli $(a,w)$ jest styczny,
to $ta=z\cdot w$ i
\[
 |w|^2-a^2
 =|w|^2-\frac{(z\cdot w)^2}{t^2}
 \geq\frac{|w|^2}{t^2}>0
 \quad\text{dla }(a,w)\ne0.
\]
Możemy także opisać hiperboloidę \emph{samymi
współrzędnymi przestrzennymi} $z\in\R^n$:
\begin{equation}\label{eq:hiperboloida-wspolrzedne-metryka}
 Q(z)=(\sqrt{1+|z|^2},z),\qquad
 (g_{\mathcal H})_{ij}
   =\delta_{ij}-\frac{z_i z_j}{1+|z|^2}.
\end{equation}
Istotnie, dla $t(z)=\sqrt{1+|z|^2}$ mamy
$\partial_i t=z_i/t$, toteż ograniczenie
$-\dd t^2+|\dd z|^2$ do wykresu $Q$ daje
\eqref{eq:hiperboloida-wspolrzedne-metryka}.
Na wektorach prostopadłych do $z$ macierz $G$
ma wartość własną $1$, a w kierunku $z$
wartość $1/(1+|z|^2)$. Stąd, również przy $z=0$
przez ciągłość,
\begin{equation}\label{eq:hiperboloida-macierz-odwrotna}
 G^{-1}=I+zz^{\mathsf T},\qquad
 \det G=\frac1{1+|z|^2}.
\end{equation}
Równość dla macierzy odwrotnej można sprawdzić
mnożąc $(I-zz^{\mathsf T}/t^2)(I+zz^{\mathsf T})$
i używając $t^2=1+|z|^2$.

\emph{Dlaczego to ta sama geometria?}
Odwzorowanie
\[
 P:B^n\longrightarrow\mathcal H^n,\qquad
 P(u)=\left(\frac{1+|u|^2}{1-|u|^2},
            \frac{2u}{1-|u|^2}\right)
\]
ma odwrotność $u=z/(t+1)$.
Dla $q=1-|u|^2$ różniczkujemy:
\[
 \dd t=\frac{4(u\cdot\dd u)}{q^2},
 \qquad
 \dd z=\frac{2\,\dd u}{q}
          +\frac{4u(u\cdot\dd u)}{q^2}.
\]
Po podniesieniu do kwadratu w
$|\,\dd z\,|^2-\dd t^2$ współczynnik przy
$(u\cdot\dd u)^2$ jest
$16(q+|u|^2-1)/q^4=0$. Pozostaje dokładnie
$4|\,\dd u\,|^2/q^2$, czyli $P^*g_{\mathcal H}=g_B$.

Kula przechodzi na półprzestrzeń przez
\[
 F(u)=\left(\frac{2v}{D},\frac{1-|u|^2}{D}\right),
 \quad u=(v,s),\quad D=|v|^2+(s-1)^2.
\]
Jest to dyfeomorfizm: dla $(x,y)$ z $y>0$
odwrotność ma postać
\[
 F^{-1}(x,y)=
 \left(\frac{2x}{E},
       \frac{|x|^2+y^2-1}{E}\right),
 \quad E=|x|^2+(y+1)^2.
\]
Aby sprawdzić metrykę bez długiego rozwijania,
zauważmy, że $F$ jest złożeniem przesunięcia,
inwersji $w\mapsto w/|w|^2$, jednokładności
o skali $2$ i odbicia. Różniczka inwersji to
$|w|^{-2}(I-2ww^{\mathsf T}/|w|^2)$;
macierz w nawiasie jest ortogonalna.
Zatem $|\,\dd F_u(h)\,|=2|h|/D$.
Współrzędna wysokości wynosi
$y=(1-|u|^2)/D$, więc
\[
 F^*g_{\mathbb H}(h,h)
 =\frac{4|h|^2/D^2}{(1-|u|^2)^2/D^2}
 =g_B(h,h).
\]
Oba odwzorowania zachowują więc długości
wektorów i krzywych. Sfera i przestrzeń
hiperboliczna są zakrzywione inaczej;
formalny opis ich krzywizny odłożymy do
rozdziału o tensorze krzywizny.
\end{przyklad}

\begin{figure}[H]
\centering
\begin{tikzpicture}[>=Latex,scale=1]
  \begin{scope}[shift={(-4.25,0)}]
    \fill[manifoldblue!7] (-1.3,-1.03) rectangle (1.3,1.10);
    \draw[manifoldblue!70!black,thick] (-1.3,-1.03)--(1.3,-1.03)
      node[right] {$y=0$};
    \fill[manifoldred] (-1.04,-.68) circle (2pt);
    \draw[manifoldred,very thick,->] (-1.04,-.68)--(-.69,-.68);
    \fill[manifoldred] (-1.04,.36) circle (2pt);
    \draw[manifoldred,very thick,->] (-1.04,.36)--(.35,.36);
    \node[manifoldblue!75!black] at (0,-1.48)
       {półprzestrzeń};
  \end{scope}
  \begin{scope}[shift={(0,0)}]
    \fill[manifoldblue!7] (0,0) circle (1.15);
    \draw[manifoldblue!70!black,very thick] (0,0) circle (1.15);
    \fill[manifoldred] (-.47,-.28) circle (2pt);
    \pgfmathsetmacro{\vlen}{1.15*(1-(.47*.47+.28*.28)/(1.15*1.15))/2}
    \draw[manifoldred,very thick,->] (-.47,-.28)--({-.47+\vlen},-.28);
    \fill[manifoldred] (.76,.35) circle (2pt);
    \pgfmathsetmacro{\vlen}{1.15*(1-(.76*.76+.35*.35)/(1.15*1.15))/2}
    \draw[manifoldred,very thick,->] (.76,.35)--({.76+\vlen},.35);
    \node[manifoldblue!75!black] at (0,-1.48)
       {kula Poincarégo};
  \end{scope}
  \begin{scope}[shift={(4.25,-.05)}]
    \pgfmathsetmacro{\htop}{3.4*(sqrt(1+1.22*1.22)-1)-1.03}
    \shade[top color=manifoldblue!18,bottom color=manifoldblue!47]
      plot[domain=-1.22:1.22,samples=61]
        (\x,{3.4*(sqrt(1+\x*\x)-1)-1.03})
      arc[start angle=0,end angle=180,x radius=1.22,y radius=.30]
      --cycle;
    \draw[manifoldblue!75!black,thick]
      plot[domain=-1.22:1.22,samples=61]
        (\x,{3.4*(sqrt(1+\x*\x)-1)-1.03});
    \draw[manifoldblue!75!black,thick] (0,\htop)
      ellipse (1.22 and .30);
    \foreach \rad in {.50,.88}{
      \pgfmathsetmacro{\height}{3.4*(sqrt(1+\rad*\rad)-1)-1.03}
      \draw[manifoldblue!50!black,thin] (0,\height)
        ellipse ({\rad} and {.246*\rad});
    }
    \node[manifoldblue!75!black] at (0,-1.43)
       {hiperboloida};
  \end{scope}
\end{tikzpicture}
\caption{Trzy zapisy tej samej geometrii hiperbolicznej. Czerwone strzałki przedstawiają wektory o długości hiperbolicznej $1$: taki wektor ma różną długość na euklidesowym rysunku półprzestrzeni i kuli. Linia $y=0$ oraz okrąg brzegowy nie należą do rozmaitości. Rysunek hiperboloidy przedstawia fragment górnego płata; jego profil jest gładki także w najniższym punkcie.}
\label{fig:metryki-trzy-modele-hiperboliczne}
\end{figure}

\section{Metryki Schwarzschilda i Kerra: inne znaki}
\label{sec:metryki-lorentzowskie}

\begin{definicja}[Metryka lorentzowska]
\index{metryka lorentzowska@Metryka lorentzowska}
\label{def:metryka-lorentzowska}
Metryka pseudoriemannowska jest gładkim,
symetrycznym, \emph{niezdegenerowanym} tensorem
typu $(0,2)$; dopuszczamy ujemne wartości $g(v,v)$.
W czterech wymiarach metryka o jednym kierunku
ujemnym i trzech dodatnich ma sygnaturę $(-,+,+,+)$
i nazywa się \emph{lorentzowską}.
\end{definicja}

Tu słowo „metryka” oznacza tensor, a nie funkcję
odległości na parach punktów. Dla wektora
czasopodobnego $v$ może być $g(v,v)<0$, a dla
niezerowego wektora światłopodobnego $g(v,v)=0$.
Wzory na riemannowską długość i odległość
z następnej sekcji nie stosują się do całej
takiej czasoprzestrzeni.

W poniższych dwóch przykładach używamy jednostek
geometrycznych $G=c=1$. W jednostkach fizycznych
parametr masy ma wymiar długości:
$M=G M_{\mathrm{fiz}}/c^2$, gdzie
$M_{\mathrm{fiz}}$ jest masą ciała, $G$ stałą
grawitacji, a $c$ prędkością światła. Współrzędna
$t$ ma wtedy wymiar długości ($t=c\,t_{\mathrm{fiz}}$).
Kąty $\theta\in(0,\pi)$ i $\phi\pmod{2\pi}$
oznaczają odpowiednio kąt od osi obrotu oraz azymut.
Symbole $\dd t^2$, $\dd r^2$ itd. oznaczają
kwadraty różniczek współrzędnych w zapisie
symetrycznego tensora $g$.

\begin{przyklad}[Schwarzschild]
\label{prz:metryka-schwarzschild}
W jednostkach geometrycznych $G=c=1$, dla $M>0$
i na obszarze zewnętrznym $r>2M$, metryka ma postać
\begin{equation}\label{eq:schwarzschild}
 g_{\mathrm{Sch}}
 =-\left(1-\frac{2M}{r}\right)\dd t^2
  +\left(1-\frac{2M}{r}\right)^{-1}\dd r^2
  +r^2\bigl(\dd\theta^2+\sin^2\theta\,\dd\phi^2\bigr).
\end{equation}
Parametr $M$ opisuje masę źródła, natomiast $r$
jest \emph{promieniem powierzchniowym}:
sfera określona przez $t,r=\mathrm{const}$
ma pole $4\pi r^2$. Nie jest to ogólnie odległość
zmierzona wzdłuż promienia; z metryki na
$t=\mathrm{const}$ odczytujemy element takiej
odległości $\dd\ell=\dd r/\sqrt{1-2M/r}$.
Współrzędna $t$ jest unormowana tak, aby asymptotycznie, gdy
$r\to\infty$, zgadzać się z czasem własnym obserwatorów
statycznych po przeliczeniu jednostek. Dla obserwatora
o stałych $r,\theta,\phi$ odczytujemy z metryki
$\dd\tau=\sqrt{1-2M/r}\,\dd t$; przy skończonym $r$
czas własny $\tau$ różni się więc od $t$.
Przy $r>2M$ współczynnik $\dd t^2$ jest ujemny,
a trzy przestrzenne są dodatnie, więc sygnatura
wynosi $(-,+,+,+)$. Dla ustalonego $t$
otrzymujemy \emph{riemannowską} metrykę przestrzennego
przekroju
\[
 h_{\mathrm{Sch}}
 =\left(1-\frac{2M}{r}\right)^{-1}\dd r^2
  +r^2\bigl(\dd\theta^2+\sin^2\theta\,\dd\phi^2\bigr).
\]
Na sferze $t,r=\mathrm{const}$ metryka indukowana
to \eqref{eq:metryka-sfery} z $R=r$, więc jej
pole wynosi $4\pi r^2$. Współrzędne sferyczne
nie działają na biegunach, a postać
\eqref{eq:schwarzschild} nie obejmuje $r=2M$;
odpowiednie przedłużenie opisuje się innymi
współrzędnymi.
\end{przyklad}

\begin{przyklad}[Kerr]
\label{prz:metryka-kerr}
Przy $M>0$, $|a|\leq M$ i
$r>r_+:=M+\sqrt{M^2-a^2}$ połóżmy
$\Sigma=r^2+a^2\cos^2\theta$ i
$\Delta=r^2-2Mr+a^2$. W zewnętrznej
dziedzinie współrzędnych Boyera--Lindquista
metrykę zapisuje się zwięźle jako
\begin{align}\label{eq:kerr}
 g_{\mathrm K}
 ={}&-\frac{\Delta}{\Sigma}
       (\dd t-a\sin^2\theta\,\dd\phi)^2
  +\frac{\sin^2\theta}{\Sigma}
       \bigl((r^2+a^2)\dd\phi-a\,\dd t\bigr)^2\\
 &+\frac{\Sigma}{\Delta}\dd r^2
  +\Sigma\,\dd\theta^2.\notag
\end{align}
Tutaj $M$ ponownie jest parametrem masy,
$J$ oznacza moment pędu źródła, a
$a=J/M$ w jednostkach $G=c=1$ jest parametrem
\emph{obrotu} o wymiarze długości
(w zwykłych jednostkach $a=J_{\mathrm{fiz}}/
(M_{\mathrm{fiz}}c)$). Znak $a$ ustala zwrot
obrotu względem rosnącego $\phi$; $a=0$ daje
przypadek bez obrotu. Założenie $|a|\leq M$
zapewnia rzeczywiste pierwiastki $\Delta=0$;
$r_+$ jest większym z nich, czyli współrzędną
zewnętrznego horyzontu. Wielkości $\Sigma$ i
$\Delta$ są skrótami zależnymi od współrzędnych,
nie dodatkowymi parametrami. Współrzędna $r$
nie mierzy tu odległości przestrzennej i przy
$a\ne0$ nie jest także promieniem sfery o polu
$4\pi r^2$; $t,\theta,\phi$ są współrzędnymi
Boyera--Lindquista. Wzór opisuje zewnętrze
horyzontu, a nie całą czasoprzestrzeń.
Kwadrat nawiasu jest kwadratem symetrycznym
$1$-formy. Ponieważ $\Delta>0$ dla $r>r_+$
i $0<\theta<\pi$, cztery wyświetlone
kwadraty mają odpowiednio znaki $-,+,+,+$
(po liniowo niezależnej zamianie pary
$\dd t,\dd\phi$). Istotnie, macierz tej zamiany to
$\left(\begin{smallmatrix}1&-a\sin^2\theta\\-a&r^2+a^2\end{smallmatrix}\right)$,
a jej wyznacznik wynosi $r^2+a^2\cos^2\theta=\Sigma>0$. W szczególności nie jest
to metryka riemannowska.

Sprawdźmy związek ze Schwarzschildem:
dla $a=0$ mamy $\Sigma=r^2$ oraz
$\Delta=r^2(1-2M/r)$; cztery składniki
\eqref{eq:kerr} stają się dokładnie
składnikami \eqref{eq:schwarzschild}.
Dla $a\ne0$ rozwinięcie dwóch pierwszych
kwadratów daje wyraz mieszany
\[
 -\frac{4Mar\sin^2\theta}{\Sigma}
      \,\dd t\,\dd\phi.
\]
W macierzy tensora oznacza to
$g_{t\phi}=-2Mar\sin^2\theta/\Sigma$:
oba wyrazy $g_{t\phi}\dd t\otimes\dd\phi$
i $g_{\phi t}\dd\phi\otimes\dd t$ składają
się na współczynnik podwojony.
Ten wyraz wiąże czas z kierunkiem obrotu.
Przekrój $t=\mathrm{const}$ ma dodatnio określoną
metrykę w podanej dziedzinie. Współczynniki przy
$\dd r^2$ i $\dd\theta^2$ są dodatnie, a współczynnik
przy $\dd\phi^2$ wynosi
\[
 \frac{\sin^2\theta}{\Sigma}
 \bigl((r^2+a^2)^2-a^2\Delta\sin^2\theta\bigr)>0,
\]
ponieważ wyrażenie w nawiasie równa się
$(r^2+a^2)\Sigma+2Mr a^2\sin^2\theta$.
Pomiary przestrzenne zależą jednak od wybranego
przekroju czasoprzestrzeni.
\end{przyklad}

W obu przykładach zapisaliśmy konkretne tensory
i sprawdziliśmy ich znaki oraz prostą zależność
między nimi. Ich wyprowadzenie z równań Einsteina
wymaga aparatu krzywizny, którego jeszcze nie
wprowadziliśmy.

\section{Długość, energia i odległość}
\label{sec:metryka-dlugosc-energia}

\begin{definicja}[Długość i energia krzywej]
\index{dlugosc i energia krzywej@Długość i energia krzywej}
\label{def:dlugosc-energia}
Niech $a<b$, a $\gamma:[a,b]\to(M,g)$ będzie ciągła
i odcinkami klasy $C^1$ (na skończonym podziale przedziału).
Wektor prędkości to $\dot\gamma(t)$, a jego długość,
nazywana także szybkością, wynosi $|\dot\gamma(t)|_g$.
Jej \emph{długość} i
\emph{energia} to odpowiednio
\begin{equation}\label{eq:dlugosc-energia}
 L_g(\gamma)=\int_a^b|\dot\gamma(t)|_g\,\dd t,
 \qquad
 E_g(\gamma)=\frac12\int_a^b
                  |\dot\gamma(t)|_g^2\,\dd t.
\end{equation}
Dla krzywych łączących $p$ z $q$ definiujemy
\begin{equation}\label{eq:odleglosc-riemannowska}
 d_g(p,q)=\inf_{\gamma(a)=p,\ \gamma(b)=q}
                 L_g(\gamma).
\end{equation}
Infimum bierzemy po krzywych odcinkami $C^1$.
Jeśli nie ma takiej krzywej, przyjmujemy $\inf\varnothing=+\infty$.
Zatem dla rozmaitości niespójnej wzór daje metrykę
skończoną osobno na każdej składowej; między różnymi
składowymi ma wartość $+\infty$.
\end{definicja}

Metryka $g$ mierzy wektor prędkości w każdym
punkcie, długość sumuje te pomiary wzdłuż
\emph{jednej} drogi, a $d_g$ szuka najkrótszej
możliwej długości. Infimum nie zakłada jeszcze,
że istnieje krzywa, która je osiąga.

\begin{stwierdzenie}[Zmiana parametru i energia]
\label{prop:energia-dlugosc}
Jeśli $h:[c,d]\to[a,b]$ jest rosnącym dyfeomorfizmem
klasy $C^1$, to
$L_g(\gamma\circ h)=L_g(\gamma)$.
Ponadto
\begin{equation}\label{eq:energia-cauchy}
 L_g(\gamma)^2\leq 2(b-a)E_g(\gamma),
\end{equation}
z równością wtedy i tylko wtedy, gdy prędkość
$|\dot\gamma|_g$ jest stała prawie wszędzie.
Dla regularnej krzywej klasy $C^1$, czyli takiej, że
$\dot\gamma(t)\ne0$, istnieje przeparametryzowanie klasy
$C^1$ ze stałą szybkością. Osiąga ono najmniejszą energię
wśród wszystkich regularnych przeparametryzowań na
ustalonym $[a,b]$. Dla krzywej odcinkami regularnej
odpowiednie przeparametryzowanie jest odcinkami klasy $C^1$.
\end{stwierdzenie}
\begin{proof}
Reguła łańcuchowa daje
$(\gamma\circ h)'(s)=\dot\gamma(h(s))h'(s)$.
Ponieważ $h'>0$, jednorodność normy i podstawienie
$t=h(s)$ dają
\[
 L_g(\gamma\circ h)
 =\int_c^d|\dot\gamma(h(s))|_g h'(s)\,\dd s
 =\int_a^b|\dot\gamma(t)|_g\,\dd t.
\]
Dla malejącego dyfeomorfizmu ten sam rachunek używa
$|h'|=-h'$ i odwróconych granic całkowania, więc długość
także się nie zmienia. Dla energii w pierwszej całce występuje
$(h')^2$, a więc podstawienie nie usuwa całego
czynnika: energia na ogół się zmienia.
Nierówność Cauchy'ego--Schwarza dla funkcji
$1$ i $|\dot\gamma|_g$ na $[a,b]$ daje
\[
 \left(\int_a^b|\dot\gamma|_g\,\dd t\right)^2
 \leq\left(\int_a^b1\,\dd t\right)
      \left(\int_a^b|\dot\gamma|_g^2\,\dd t\right).
\]
To \eqref{eq:energia-cauchy}. Warunek równości
w Cauchym--Schwarzu oznacza, że obie funkcje
są proporcjonalne prawie wszędzie, czyli
$|\dot\gamma|_g$ jest stała prawie wszędzie.
Dla regularnej krzywej $|\dot\gamma|_g>0$;
funkcja
$s(t)=a+\frac{b-a}{L_g(\gamma)}
               \int_a^t|\dot\gamma(\tau)|_g\,\dd\tau$
jest rosnąca, przenosi końce przedziału na
siebie, a reguła łańcuchowa pokazuje,
że $\gamma\circ s^{-1}$ ma stałą prędkość
$L_g(\gamma)/(b-a)$. Osiąga zatem dolną granicę
energii $L_g(\gamma)^2/[2(b-a)]$.
\end{proof}

\begin{twierdzenie}[Odległość riemannowska jest metryką]
\label{thm:odleglosc-jest-metryka}
Jeżeli $M$ jest spójna, to \eqref{eq:odleglosc-riemannowska}
przyjmuje skończone wartości, spełnia aksjomaty
metryki przestrzeni metrycznej $(M,d_g)$
i wyznacza pierwotną topologię rozmaitości.
\end{twierdzenie}
\begin{proof}
W wymiarze $0$ spójna rozmaitość składa się z jednego
punktu i teza jest natychmiastowa. Załóżmy $n\geq1$.
Relacja „można połączyć krzywą odcinkami $C^1$” jest
relacją równoważności. Każda jej klasa jest otwarta:
w małej kuli współrzędnych, a przy brzegu w półkuli,
można połączyć dwa punkty odcinkiem prostym w mapie.
Dopełnienie klasy także jest otwarte, jako suma innych
klas. Ze spójności wynika, że jest tylko jedna klasa.
Na każdym ze skończenie wielu kawałków takiej krzywej
prędkość i współczynniki metryki są ciągłe na zwartym
przedziale, więc długość jest skończona.
Zatem $d_g(p,q)<\infty$.

Długości są nieujemne, stała krzywa ma
długość zero, odwrócenie kierunku zachowuje
długość, a długości sklejonych krzywych
dodają się. Biorąc infima (i drogi o długości
dowolnie bliskiej obu infimom), otrzymujemy
$d_g(p,p)=0$, symetrię i nierówność trójkąta.

Zostało najważniejsze: $d_g(p,q)>0$ dla $p\ne q$.
Wybierzmy mapę $x$ z $x(p)=0$ i domkniętą
kulę współrzędnych $\overline B_{2r}$ leżącą
w obrazie tej mapy. Jeśli $p\in\partial M$, wybieramy
mapę przybrzegową z $x(p)=0$ i w całym poniższym
rachunku zastępujemy $B_\rho$ przez
$B_\rho\cap\{x^n\geq0\}$. Te półkule są wypukłe,
więc odcinki używane w dowodzie pozostają w $M$.
„Wyjście z kuli” oznacza wtedy przekroczenie jej
sferycznej części brzegu wewnątrz półprzestrzeni.
Na zwartym zbiorze
$x^{-1}(\overline B_{2r})$ dodatnia określoność
i ciągłość macierzy $G$ dają stałe $c,C>0$,
takie że
\[
 c|w|\leq\sqrt{w^{\mathsf T}G(x)w}\leq C|w|
 \quad (x\in\overline B_{2r}).
\]
Istnienie wspólnego $c$ można zobaczyć,
minimalizując $w^{\mathsf T}G(x)w$ po zwartym
iloczynie $\overline B_{2r}\times S^{n-1}$.

Jeśli $q$ leży w $x^{-1}(B_r)$, droga z $p$
do $q$ pozostająca w $B_{2r}$ ma długość
co najmniej $c|x(q)|$, bo długość euklidesowa
drogi jest nie mniejsza niż przemieszczenie
końców. Droga wychodząca z $B_{2r}$ ma do
pierwszego wyjścia długość co najmniej $2cr$,
więc również co najmniej $c|x(q)|$.
Jeśli $q$ leży poza $x^{-1}(B_r)$, każda droga
do $q$ musi najpierw opuścić $B_r$, zatem
ma długość co najmniej $cr$. Oba przypadki
dają dodatnią odległość od $p$.

Te same oszacowania porównują topologie.
Dla $q$ w $B_r$ odcinek prosty od $0$ do $x(q)$
daje $d_g(p,q)\leq C|x(q)|$.
Z drugiej strony kula metryczna
$B_{d_g}(p,cr)$ leży w $x^{-1}(B_r)$,
a na niej $d_g(p,q)\geq c|x(q)|$.
Zbieżność w mapie i zbieżność dla $d_g$
są więc lokalnie równoważne.
\end{proof}

\begin{przyklad}[Niezupełność i brzeg modelu hiperbolicznego]
Przedział $(0,1)$ z $g=\dd x^2$ ma zwykłą
odległość $|p-q|$, ale ciąg $1/j$, $j\geq2$, jest
Cauchy'ego i nie zbiega w tym przedziale.
Samo istnienie $g$ nie zapewnia zupełności.

W półprzestrzeni hiperbolicznej pionowa droga
od $(x,y_0)$ do $(x,y_1)$ ma długość
$|\log(y_1/y_0)|$. Każda inna droga ma
\[
 L_g(\gamma)\geq\int\frac{|y'(t)|}{y(t)}\,\dd t
 \geq\left|\int\frac{y'(t)}{y(t)}\,\dd t\right|
 =|\log(y_1/y_0)|,
\]
więc jest to rzeczywista odległość tych punktów.
Granica $y=0$ leży nieskończenie daleko.
Podobnie w kuli Poincarégo
$d_g(0,u)=\log\frac{1+|u|}{1-|u|}$,
bo na każdej drodze $\gamma$ funkcja
$\rho(t)=|\gamma(t)|$ jest absolutnie ciągła
i $|\rho'|\leq|\dot\gamma|$ prawie wszędzie.
Stąd
\[
 L_g(\gamma)\geq\int\frac{2|\rho'|}{1-\rho^2}\,\dd t
 \geq\left|\int\frac{2\rho'}{1-\rho^2}\,\dd t\right|
 =\int_0^{|u|}\frac{2\,\dd s}{1-s^2}
 =\log\frac{1+|u|}{1-|u|}.
\]
Promień od $0$ do $u$ osiąga tę dolną granicę.
\end{przyklad}

\section{Objętość wyznaczona przez metrykę}
\label{sec:metryka-objetosc}

Jeden wektor ma długość; $n$ niezależnych
wektorów stycznych rozpinających maleńki
równoległościan ma objętość. W bazie
ortogonalnej jest to iloczyn długości.
W dowolnej bazie odpowiednią liczbę
zawiera wyznacznik macierzy Grama.

\begin{definicja}[Element objętości i miara riemannowska]
\index{element objetosci i miara riemannowska@Element objętości i miara riemannowska}
\label{def:miara-riemannowska}
W mapie $x$ na $(M^n,g)$ \emph{element objętości}
jest dodatnią gęstością
\begin{equation}\label{eq:element-objetosci}
 \dd\mu_g=\sqrt{\det(g_{ij}(x))}\,
                 |\,\dd x^1\cdots\dd x^n\,|.
\end{equation}
Przez całkowanie tej gęstości w mapach
otrzymujemy miarę borelowską $\mu_g$.
Dla funkcji borelowskiej $f\geq0$ lub funkcji całkowalnej,
o nośniku zawartym w jednej mapie,
\[
 \int_M f\,\dd\mu_g
  =\int_{x(U)} f(x^{-1}(u))
      \sqrt{\det G(u)}\,\dd u^1\cdots\dd u^n.
\]
Ogólną całkę definiujemy z podziałem jedności:
dla funkcji nieujemnej dopuszczamy wartość $+\infty$,
a dla funkcji ze znakiem wymagamy $\int_M|f|\,\dd\mu_g<\infty$.
Sam zwarty nośnik funkcji mierzalnej nie zapewnia całkowalności.
W wymiarze $0$ przyjmujemy wyznacznik macierzy pustej równy $1$;
$\mu_g$ jest wtedy miarą liczącą punkty. Na rozmaitości zorientowanej
\emph{forma objętościowa} ma w dodatniej mapie postać
$\mathrm{vol}_g=\sqrt{\det G}\,
 \dd x^1\wedge\cdots\wedge\dd x^n$.
Dla $n=0$ forma objętościowa ma w punkcie wartość
$\varepsilon_p\in\{1,-1\}$ zadaną jego orientacją;
całkowanie z tą samą orientacją daje dodatnią masę $1$.
\end{definicja}

Kreski pionowe w \eqref{eq:element-objetosci}
oznaczają gęstość: zmiana orientacji nie zmienia
miary. Forma $\mathrm{vol}_g$ ma natomiast znak
zależny od orientacji. Dlatego pole powierzchni
możemy obliczyć również na rozmaitości
nieorientowalnej, choć nie ma na niej globalnej
formy objętościowej.

\begin{twierdzenie}[Niezależność elementu objętości od mapy]
\label{thm:element-objetosci-globalny}
Wzór \eqref{eq:element-objetosci} daje
jednoznaczną miarę na $M$, niezależną od atlasu
i podziału jedności. Jeżeli $F:(M,g)\to(N,h)$
jest izometrią, czyli dyfeomorfizmem z $F^*h=g$,
to $\mu_g(F^{-1}A)=\mu_h(A)$ dla zbiorów
borelowskich $A\subset N$.
\end{twierdzenie}
\begin{proof}
Na przecięciu map niech $x=x(y)$ i
$J=\partial x/\partial y$. Ze
Stwierdzenia~\ref{prop:metryka-zmiana} wynika
\[
 \sqrt{\det G_y}
   =|\det J|\sqrt{\det G_x}.
\]
To dokładnie czynnik w twierdzeniu o zmianie
zmiennych w całce Lebesgue'a, więc obliczenia
tego samego zbioru w obu mapach są zgodne.
Wybierzmy przeliczalne pokrycie mapami
$(U_j)$, a mierzalny zbiór $A$ rozbijmy
na rozłączne fragmenty
$A\cap U_1$,
$A\cap(U_2\setminus U_1)$ itd.
Sumujemy ich lokalne miary. Na przecięciach
otrzymalibyśmy te same liczby w innych mapach,
więc wynik nie zależy od tego rozbicia;
lokalna addytywność całki daje przeliczalną
addytywność globalną. Analogicznie porównujemy
dwa podziały jedności przez ich wspólne
rozdrobnienie.

Dla izometrii macierz jej różniczki spełnia
w lokalnych bazach $G_M=(DF)^{\mathsf T}G_N(DF)$.
Poprzedni rachunek wyznaczników i podstawienie
w całce pokazują
$\int_M(f\circ F)\dd\mu_g=\int_N f\dd\mu_h$
dla nieujemnych funkcji mierzalnych $f$.
Wstawiając funkcję charakterystyczną $A$,
otrzymujemy tezę o miarach.
\end{proof}

\begin{przyklad}[Gęstości w modelach]
Ze \eqref{eq:metryka-sfery} otrzymujemy
$\dd\mu_{S_R^2}=R^2\sin\theta\,
\dd\theta\,\dd\phi$ dla $0<\theta<\pi$,
a stąd
$\operatorname{Pole}(S_R^2)
=\int_0^{2\pi}\int_0^\pi
R^2\sin\theta\,\dd\theta\,\dd\phi=4\pi R^2$.
Parametry $0<\theta<\pi$, $0<\phi<2\pi$ pomijają
bieguny i jeden południk. Zbiory te mają zerową miarę
powierzchniową: w regularnych mapach są punktami lub
wykresami jednowymiarowymi, zerowymi dla dwuwymiarowej
miary Lebesgue'a z twierdzenia Fubiniego. Gładka gęstość
nie zmienia tej zerowości. Nie całkujemy więc przez
osobliwość mapy; dodajemy jedynie zbiór miary zero.
W półprzestrzeni hiperbolicznej macierz
metryki jest $y^{-2}I_n$, więc
$\dd\mu_{\mathbb H}=y^{-n}\dd x^1\cdots
\dd x^{n-1}\dd y$.
W kuli Poincarégo gęstość wynosi
$(2/(1-|u|^2))^n\dd u^1\cdots\dd u^n$.
W przestrzennych współrzędnych $z$ hiperboloidy
z~\eqref{eq:hiperboloida-wspolrzedne-metryka}
wyznacznik obliczony w
\eqref{eq:hiperboloida-macierz-odwrotna} daje
\[
 \dd\mu_{\mathcal H}
  =\frac{\dd z^1\cdots\dd z^n}{\sqrt{1+|z|^2}}.
\]
Podana formuła dotyczy \emph{metryki indukowanej
na hiperboloidzie}, nie miary otaczającej
przestrzeni Minkowskiego. Izometrie $P$ i $F$
z Przykładu~\ref{prz:trzy-modele-hiperboliczne}
przekształcają te trzy gęstości w siebie.

Dla przestrzennego przekroju Schwarzschilda
z Przykładu~\ref{prz:metryka-schwarzschild}
macierz $h_{\mathrm{Sch}}$ jest diagonalna:
\[
 \det h_{\mathrm{Sch}}
  =\frac{r^4\sin^2\theta}{1-2M/r},
 \qquad
 \dd\mu_{h_{\mathrm{Sch}}}
  =\frac{r^2\sin\theta}{\sqrt{1-2M/r}}\,
       \dd r\,\dd\theta\,\dd\phi.
\]
Jest to \emph{trójwymiarowa} miara wybranego
przekroju, a nie miara całej metryki lorentzowskiej.
\end{przyklad}

\section{Metryka indukowana i całki po podrozmaitości}
\label{sec:metryka-indukowana}

\begin{definicja}[Metryka indukowana]
\index{metryka indukowana@Metryka indukowana}
\label{def:metryka-indukowana}
Niech $i:S^k\hookrightarrow(M^n,g)$ będzie
gładką podrozmaitością, $k\leq n$.
Na $S$ definiujemy $h=i^*g$:
\[
 h_p(v,w)=g_{i(p)}(\dd i_pv,\dd i_pw)
 \quad(v,w\in T_pS).
\]
Ponieważ $\dd i_p$ jest iniektywna, $h$ jest
metryką riemannowską. Miara $\mu_h$ na $S$
pozwala całkować funkcję $f:S\to\R$;
oznaczenie $\int_S f\,\dd A_g$ stosujemy,
gdy $k=n-1$.
\end{definicja}

Parametryzacja $F:U\subset\R^k\to S$
wybiera wektory $F_1,\ldots,F_k$ styczne
do powierzchni. Macierz Grama tych wektorów
wyznacza rzeczywisty rozmiar
obrazu małego prostokąta współrzędnych.

\begin{figure}[!htbp]
\centering
\begin{tikzpicture}[>=Stealth,line cap=round,line join=round]
 \fill[manifoldgold!11]
   (0,0)--(4.0,0)--(4.85,1.05)--(.85,1.05)--cycle;
 \draw[manifoldgold!70!black,thick]
   (0,0)--(4.0,0)--(4.85,1.05)--(.85,1.05)--cycle;
 \fill[manifoldblue!13]
   (0,1.15) .. controls (1.4,1.45) and (2.8,1.18) .. (4,1.68)
   --(4.85,2.85)
   .. controls (3.4,2.36) and (2.0,2.65) .. (.85,2.18)--cycle;
 \draw[manifoldblue!75!black,thick]
   (0,1.15) .. controls (1.4,1.45) and (2.8,1.18) .. (4,1.68)
   --(4.85,2.85)
   .. controls (3.4,2.36) and (2.0,2.65) .. (.85,2.18)--cycle;
 \draw[manifoldblue!55,dashed]
   (.85,2.18)--(.85,1.05)
   (4,1.68)--(4,0)
   (4.85,2.85)--(4.85,1.05)
   (0,1.15)--(0,0);
 \coordinate (p) at (2.32,1.77);
 \fill[manifoldred] (p) circle (2.6pt);
 \draw[->,line width=1.2pt,manifoldblue!80!black]
   (p)--(3.39,1.81) node[right,font=\small] {$F_u$};
 \draw[->,line width=1.2pt,manifoldgreen!80!black]
   (p)--(2.79,2.40) node[above,font=\small] {$F_v$};
 \draw[->,line width=1.2pt,manifoldred!80!black]
   (p)--(2.08,3.11) node[above,font=\small] {$\nu$};
 \node[manifoldgold!70!black,font=\small] at (2.10,.30) {$D\subset\R^2$};
 \node[manifoldblue!75!black,font=\small] at (4.03,2.25) {$F(D)$};
\end{tikzpicture}
\caption{Parametryzacja unosi obszar parametrów $D$
na powierzchnię. Wektory $F_u,F_v$ wyznaczają
element pola $\sqrt{\det(g(F_a,F_b))}\,\dd u\,\dd v$;
normalna $\nu$ pozwala dodatkowo mierzyć strumień.}
\label{fig:metryka-indukowana-pole}
\end{figure}

\begin{twierdzenie}[Wzór parametryczny na całkę]
\label{thm:calka-podrozmaitosc}
Jeżeli $F$ jest parametryzacją kawałka
podrozmaitości, jeden raz pokrywającą
rozważany fragment, i $f$ jest całkowalna
o nośniku w tym fragmencie, to
\begin{equation}\label{eq:calka-podrozmaitosc}
 \int_{F(U)}f\,\dd\mu_h
 =\int_U f(F(u))
   \sqrt{\det\bigl(g(F_a,F_b)\bigr)_{a,b=1}^k}
                      \,\dd u^1\cdots\dd u^k .
\end{equation}
Dla $M=\R^3$ z metryką euklidesową i $k=2$ czynnik pod pierwiastkiem
jest $|F_u\times F_v|^2$.
\end{twierdzenie}
\begin{proof}
Z definicji cofnięcia tensora
$h_{ab}(u)=g_{F(u)}(\partial_aF,\partial_bF)$.
Podstawmy tę macierz do
\eqref{eq:element-objetosci} w wymiarze $k$:
dostajemy dokładnie prawą stronę
\eqref{eq:calka-podrozmaitosc}.
Dlaczego jest to pole, a nie sztuczny czynnik?
Wybierzmy ortonormalną bazę w $T_{F(u)}M$
i zapiszmy kolumny $F_a$ w macierzy $J$.
Wówczas $(h_{ab})=J^{\mathsf T}J$.
Ortogonalizując kolumny metodą Grama--Schmidta,
otrzymujemy prostopadłe wektory o długościach
$\ell_1,\ldots,\ell_k$. Ortogonalna zmiana
kierunków i dodawanie wcześniejszych kolumn
do późniejszych nie zmieniają objętości
rozpiętego równoległościanu. Zatem jego
$k$-wymiarowa objętość wynosi
$\ell_1\cdots\ell_k$ i jednocześnie
$\det(J^{\mathsf T}J)=\ell_1^2\cdots\ell_k^2$.
Zmiana parametrów $u=u(v)$ mnoży ten czynnik
przez $|\det(\partial u/\partial v)|$,
jak w Stwierdzeniu~\ref{prop:metryka-zmiana};
twierdzenie o podstawieniu pokazuje
niezależność wyniku od parametryzacji.
W $\R^3$ tożsamość
$|F_u\times F_v|^2
=|F_u|^2|F_v|^2-(F_u\cdot F_v)^2$
jest wyznacznikiem macierzy Grama $2\times2$.
\end{proof}

Jeśli parametryzacja pokrywa fragment kilka
razy, prawa strona liczy go z krotnością.
Dla całej podrozmaitości używamy map
i podziału jedności, żeby każdy punkt
został policzony tylko raz. Dla $k=1$
wzór jest całką po łuku
$\int f(\gamma(t))|\dot\gamma(t)|_g\,\dd t$;
dla $k=n$ jest zwykłą całką objętościową.

\begin{przyklad}[Wykres funkcji i paraboloida]
\label{prz:calka-po-wykresie}
Dla $F(x,y)=(x,y,f(x,y))$ w euklidesowej $\R^3$
wektory styczne to $(1,0,f_x)$ i $(0,1,f_y)$,
a ich macierz Grama to
\[
 \begin{pmatrix}1+f_x^2&f_xf_y\\
                 f_xf_y&1+f_y^2\end{pmatrix}.
\]
Jej wyznacznik jest $1+f_x^2+f_y^2$, więc
\begin{equation}\label{eq:calka-wykres}
 \int_{\operatorname{graf}(f|_D)}\psi\,\dd A
 =\int_D\psi(x,y,f(x,y))
        \sqrt{1+f_x^2+f_y^2}\,\dd x\,\dd y.
\end{equation}
Pierwiastek wyraża nachylenie: skośny fragment
ma większe pole niż jego rzut na płaszczyznę.

Weźmy $f=x^2+y^2$ nad dyskiem $x^2+y^2\leq a^2$.
Współrzędne biegunowe dają pole wykresu
\[
 \int_0^{2\pi}\!\int_0^a
   \sqrt{1+4r^2}\,r\,\dd r\,\dd\phi
 =\frac{\pi}{6}\bigl((1+4a^2)^{3/2}-1\bigr).
\]
W ostatnim kroku użyliśmy podstawienia
$s=1+4r^2$, $\dd s=8r\,\dd r$.
\end{przyklad}

\begin{przyklad}[Całka skalarna po sferze]
Na $S_R^2$ obliczmy całkę z funkcji $z^2$.
Ponieważ $z=R\cos\theta$, mamy
\[
 \int_{S_R^2}z^2\,\dd A
 =R^4\int_0^{2\pi}\!\int_0^\pi
       \cos^2\theta\sin\theta\,\dd\theta\,\dd\phi
 =\frac{4\pi R^4}{3}.
\]
Wewnętrzna całka wynosi $2/3$ po podstawieniu
$s=\cos\theta$. Jest to całka \emph{skalarna};
nie potrzebujemy wybrać orientacji sfery.
\end{przyklad}

\subsection*{Całki skierowane: dlaczego pojawia się znak?}
\label{subsec:calki-skierowane}

Pole powierzchni i strumień odpowiadają na dwa różne
pytania. W całce skalarnej
$\int_S f\,\dd A$ liczymy wielkość fragmentów
powierzchni: wzór~\eqref{eq:calka-podrozmaitosc}
zawiera dodatni pierwiastek z wyznacznika Grama.
Odwrócenie parametrów nie zmienia wyniku.
W \emph{całce skierowanej} pytamy, ile pola
wektorowego przechodzi przez powierzchnię
w \emph{wybranym kierunku}. Musimy ustalić,
która strona jest dodatnia; odwrócenie
normalnej zmienia znak strumienia.

\begin{przyklad}[Całka skierowana po krzywej]
\label{prz:calka-skierowana-krzywa}
Niech $\gamma:[a,b]\to\R^3$ ma wybrany kierunek
przebiegu, a $X=(P,Q,R)$ będzie polem wektorowym.
Całka skierowana wzdłuż
krzywej (często zwana \emph{całką krzywoliniową
drugiego rodzaju}) ma wzór
\begin{equation}\label{eq:calka-skierowana-krzywa}
 \int_\gamma X\cdot\dd\mathbf r
 =\int_a^b X(\gamma(t))\cdot\dot\gamma(t)\,\dd t
 =\int_\gamma(P\,\dd x+Q\,\dd y+R\,\dd z).
\end{equation}
Jeżeli $\bar\gamma(t)=\gamma(a+b-t)$, to
$\dot{\bar\gamma}(t)=-\dot\gamma(a+b-t)$.
Podstawienie $s=a+b-t$ w całce daje
$\int_{\bar\gamma}X\cdot\dd\mathbf r
=-\int_\gamma X\cdot\dd\mathbf r$.
Natomiast $\int_\gamma f\,\dd s
=\int_a^b f(\gamma(t))|\dot\gamma(t)|\,\dd t$
nie zmienia się przy odwróceniu przebiegu:
norma usuwa znak prędkości.
\end{przyklad}

\begin{przyklad}[Całka skierowana po powierzchni]
\label{prz:calka-skierowana-powierzchnia}
Niech $\Phi:D\subset\R^2\to S\subset\R^3$
będzie regularną parametryzacją pokrywającą rozważany
fragment jednokrotnie, a orientację
$S$ wybierzmy przez normalną
$\nu=(\Phi_u\times\Phi_v)/|\Phi_u\times\Phi_v|$.
Dla pola $X$ \emph{strumień} przez $S$ wynosi
\begin{equation}\label{eq:calka-skierowana-powierzchnia}
 \int_S X\cdot\dd\mathbf S
 :=\int_S X\cdot\nu\,\dd A
 =\int_D X(\Phi(u,v))\cdot
           (\Phi_u\times\Phi_v)\,\dd u\,\dd v.
\end{equation}
To ten sam element pola co na
Rysunku~\ref{fig:metryka-indukowana-pole},
ale z kierunkiem normalnej.
Zamiana $u$ z $v$ zmienia
$\Phi_u\times\Phi_v$ na jego przeciwieństwo;
czynnik $|\Phi_u\times\Phi_v|$ w całce skalarnej
nie zmieniłby się.

Dla wykresu $\Phi(u,v)=(u,v,f(u,v))$
i stałego pola pionowego $X=(0,0,1)$ mamy
$\Phi_u\times\Phi_v=(-f_u,-f_v,1)$.
Przy normalnej ku górze
$\int_S X\cdot\dd\mathbf S=\int_D1\,\dd u\,\dd v
=\operatorname{Pole}(D)$; przy normalnej ku dołowi
wynik ma znak minus. Natomiast
$\int_S1\,\dd A
=\int_D\sqrt{1+f_u^2+f_v^2}\,\dd u\,\dd v$
jest dodatnie przy obu orientacjach.
\end{przyklad}

\begin{stwierdzenie}[Związek z całką z formy]
\label{prop:calka-skierowana-forma}
W zorientowanej $\R^3$ ze standardową
formą objętości $\mathrm{vol}=\dd x\wedge\dd y
\wedge\dd z$ polu $X=(P,Q,R)$ odpowiada $2$-forma
\[
 \beta=\iota_X\mathrm{vol}
 =P\,\dd y\wedge\dd z+
  Q\,\dd z\wedge\dd x+
  R\,\dd x\wedge\dd y.
\]
Wówczas $\int_{(S,o)}\beta
=\int_S X\cdot\nu_o\,\dd A$.
Na zorientowanej $(M^n,g)$
$\dd\mu_g=|\mathrm{vol}_{g,o}|$ w sensie gęstości,
a więc
\[
 \int_M f\,\dd\mu_g
 =\int_{(M,o)}f\,\mathrm{vol}_{g,o}.
\]
Jeśli \emph{ta sama} forma najwyższego stopnia
$\omega$ jest całkowana po $M$ z orientacją
przeciwną $-o$, to
$\int_{(M,-o)}\omega=-\int_{(M,o)}\omega$.
\end{stwierdzenie}
\begin{proof}
Z definicji cofnięcia form z
Rozdziału~\ref{ch:formy-de-rham},
\[
 (\Phi^*\beta)(\partial_u,\partial_v)
 =\mathrm{vol}\bigl(X(\Phi),\Phi_u,\Phi_v\bigr)
 =X(\Phi)\cdot(\Phi_u\times\Phi_v).
\]
Całkowanie po $D$ daje
\eqref{eq:calka-skierowana-powierzchnia}.
Dla krzywej analogicznie
$\gamma^*(P\,\dd x+Q\,\dd y+R\,\dd z)
=X(\gamma)\cdot\dot\gamma\,\dd t$.

W dodatniej mapie
$\mathrm{vol}_{g,o}=\sqrt{\det G}\,
\dd x^1\wedge\cdots\wedge\dd x^n$,
podczas gdy wzór~\eqref{eq:element-objetosci}
na $\dd\mu_g$ zawiera wartość bezwzględną
wyznacznika zmiany współrzędnych. Stąd
równość całek dla \emph{zgodnie} wybranych
orientacji i formy. Przy odwróceniu samej
orientacji dodatnie mapy stają się ujemne,
więc całka z ustalonej formy zmienia znak.
Gdybyśmy jednocześnie zastąpili
$\mathrm{vol}_{g,o}$ nową dodatnią formą
$\mathrm{vol}_{g,-o}=-\mathrm{vol}_{g,o}$,
oba znaki by się zniosły; całka z gęstości
pozostałaby ta sama. Właśnie dlatego
przy strumieniu trzeba podać \emph{stronę}
powierzchni, a przy samym polu powierzchni nie.
\end{proof}

\section{Miara Hausdorffa: pomiar bez map i pochodnych}
\label{sec:miara-hausdorffa}

Miara riemannowska wykorzystuje gładkie mapy
i wyznacznik macierzy metryki. Miara Hausdorffa
wychodzi tylko od \emph{odległości}: pokrywamy zbiór
coraz mniejszymi częściami i sumujemy odpowiednie
potęgi ich średnic. Dzięki temu ma sens również
na zbiorach, które nie są rozmaitościami.

\begin{definicja}[Miara Hausdorffa]
\index{miara hausdorffa@Miara Hausdorffa}
\label{def:miara-hausdorffa}
Niech $(X,d)$ będzie przestrzenią metryczną,
$s\geq0$, a $A\subset X$.
Przez $\operatorname{diam}_d U$ oznaczmy
$\sup\{d(x,y):x,y\in U\}$. Kładziemy
\begin{align}\label{eq:hausdorff-def}
 \mathcal H^s_{d,\delta}(A)
 &=c_s\inf\left\{
       \sum_{j=1}^{\infty}(\operatorname{diam}_d U_j)^s:
       A\subset\bigcup_j U_j,\
       \operatorname{diam}_d U_j<\delta\right\},\\
 \mathcal H^s_d(A)
 &=\lim_{\delta\downarrow0}
                  \mathcal H^s_{d,\delta}(A).
 \notag
\end{align}
Pokrycia mogą być skończone lub przeliczalne;
puste składniki pomijamy, a puste pokrycie zbioru
pustego ma koszt $0$. Przy braku dopuszczalnego
pokrycia infimum wynosi $+\infty$.
Dla $s=k$ całkowitego stosujemy normalizację
$c_k=\omega_k/2^k$, gdzie $\omega_k$
jest objętością euklidesowej kuli jednostkowej
w $\R^k$ ($\omega_0=1$). Dla innych $s$
można wziąć dowolne dodatnie $c_s$;
wartość \emph{wymiaru Hausdorffa}
$\dim_{\mathrm H}A
=\inf\{s:\mathcal H^s_d(A)=0\}$
nie zależy od tego wyboru.
\end{definicja}

Gdy $\delta$ maleje, pokryć dopuszczalnych
jest mniej, więc infimum nie maleje:
granica w \eqref{eq:hausdorff-def} istnieje,
choć może być nieskończona.
Dla $s=0$ przyjmujemy $0^0=1$ dla
niepustego zbioru w pokryciu;
$\mathcal H^0$ liczy punkty: zbiór $N$ różnych punktów
wymaga co najmniej $N$ składników dla $\delta$ mniejszego
niż ich minimalna wzajemna odległość, a pokrycie
singletonami ma koszt $N$. Dla zbioru nieskończonego
stosujemy dolne oszacowanie do dowolnie dużych
skończonych podzbiorów i dostajemy $+\infty$.
Na prostym odcinku z odległością euklidesową $\mathcal H^1$ jest
jego długością, a na gładkiej powierzchni euklidesowej
$\mathcal H^2$ jest polem.

\begin{stwierdzenie}[Dlaczego otrzymujemy miarę borelowską]
\label{prop:hausdorff-borel}
Funkcja $\mathcal H^s_d$ jest miarą zewnętrzną.
Jeśli $d(A,B):=\inf\{d(a,b):a\in A,b\in B\}>0$, to
\[
 \mathcal H^s_d(A\cup B)=\mathcal H^s_d(A)+\mathcal H^s_d(B).
\]
Wszystkie zbiory borelowskie są dla niej mierzalne,
a na nich $\mathcal H^s_d$ jest przeliczalnie addytywna.
\end{stwierdzenie}
\begin{proof}
Dla ustalonego $\delta$ monotoniczność i wartość zero
na zbiorze pustym wynikają z definicji. Pokrywając
każdy $A_i$ z kosztem co najwyżej
$\mathcal H^s_{d,\delta}(A_i)+\eta2^{-i}$ i łącząc
pokrycia, otrzymujemy przeliczalną podaddytywność.
Jeśli suma po prawej jest nieskończona, nierówność
jest automatyczna. W przeciwnym razie puszczamy
$\eta\downarrow0$. Ponadto
\[
 \mathcal H^s_{d,\delta}\Bigl(\bigcup_i A_i\Bigr)
 \leq\sum_i\mathcal H^s_{d,\delta}(A_i)
 \leq\sum_i\mathcal H^s_d(A_i).
\]
Granica przy $\delta\downarrow0$ dowodzi
podaddytywności $\mathcal H^s_d$. Dla
$\delta<d(A,B)$ żaden składnik pokrycia nie spotyka
obu zbiorów; rozdzielenie składników daje odwrotną
nierówność dla ich sumy i następnie równość w granicy.

Wykażmy mierzalność zbioru domkniętego $F$.
Oznaczmy $\mu=\mathcal H^s_d$. Kryterium Carathéodory'ego
wymaga dla każdego $T\subset X$ równości
$\mu(T)=\mu(T\cap F)+\mu(T\setminus F)$.
Jedna nierówność jest podaddytywnością. Możemy założyć
$\mu(T)<\infty$, bo dla $\mu(T)=\infty$ ta nierówność
wymusza nieskończoność prawej strony.
Niech $T_i=\{x\in T:d(x,F)\geq1/i\}$.
Addytywność dla zbiorów oddalonych daje
$\mu(T)\geq\mu(T\cap F)+\mu(T_i)$.
Żeby przejść do $T\setminus F$, rozważmy warstwy
\[
 E_j=\{x\in T:1/(j+1)\leq d(x,F)<1/j\}.
\]
Każda skończona rodzina warstw o indeksach tej samej
parzystości składa się ze zbiorów oddalonych dodatnio:
funkcja $d(\,\cdot\,,F)$ jest $1$-lipschitzowska.
Suma ich miar nie przekracza $\mu(T)$.
Zatem $\sum_j\mu(E_j)<\infty$ i jej ogony dążą do zera.
Ponieważ $F$ jest domknięty,
$T\setminus F\subset T_i\cup\bigcup_{j\geq i}E_j$,
skąd $\mu(T\setminus F)\leq\mu(T_i)+\sum_{j\geq i}\mu(E_j)$.
Puszczając $i\to\infty$, otrzymujemy żądane kryterium.
Przypadki $F=\varnothing,X$ są natychmiastowe.
Z twierdzenia Carathéodory'ego zbiory spełniające to
kryterium tworzą $\sigma$-algebrę, na której miara
zewnętrzna jest miarą. Zawiera ona zbiory domknięte,
a więc wszystkie zbiory borelowskie.
\end{proof}

\begin{lemat}[Zmiana skali odległości]
\label{lem:hausdorff-lipschitz}
Jeśli $F:(X,d_X)\to(Y,d_Y)$ jest
$L$-lipschitzowskie z $L>0$, to dla każdego $A\subset X$
i $s\geq0$
\[
 \mathcal H^s_{d_Y}(F(A))
 \leq L^s\,\mathcal H^s_{d_X}(A).
\]
Jeśli na $A$ i $F(A)$ odwzorowanie ma odwrotność
$L'$-lipschitzowską z $L'>0$, otrzymujemy również
$(L')^{-s}\mathcal H^s_{d_X}(A)
\leq\mathcal H^s_{d_Y}(F(A))$.
\end{lemat}
\begin{proof}
Obraz pokrycia $(U_j)$ zbioru $A$ pokrywa $F(A)$,
a średnica $F(U_j)$ jest co najwyżej
$L\operatorname{diam}U_j$.
Zatem pokrycia o średnicach poniżej
$\delta$ dają pokrycia obrazu o średnicach
poniżej $L\delta$, a sumy potęg średnic
są co najwyżej $L^s$ razy większe.
Bierzemy infimum i następnie granicę
$\delta\downarrow0$. Drugą nierówność
otrzymujemy przez to samo rozumowanie
zastosowane do $F^{-1}$. Można przy tym pracować
z pokryciami wewnątrz $A$: przecięcie każdego składnika
pokrycia z $A$ nie zwiększa jego średnicy. Miara
Hausdorffa podzbioru z odległością ograniczoną jest
więc taka sama jak jego miara liczona w przestrzeni
otaczającej. Założenie $L>0$ nie wyklucza map stałych,
które są lipschitzowskie z dowolną dodatnią stałą.
\end{proof}

\begin{twierdzenie}[Zgodność z miarą riemannowską]
\label{thm:hausdorff-rownosc}
Dla spójnej rozmaitości riemannowskiej $M^n$, również
z brzegiem, na $(M^n,d_g)$ z normalizacją
\eqref{eq:hausdorff-def} mamy
$\mathcal H^n_{d_g}=\mu_g$
na zbiorach borelowskich.
Jeśli $S^k\subset M$ jest gładką
podrozmaitością osadzoną, z metryką $h=i^*g$,
to $\mathcal H^k_{d_g|_S}=\mu_{i^*g}$
na $S$. Tę samą miarę na $S$ daje
odległość wewnętrzna $d_{i^*g}$, liczona na każdej
składowej $S$ osobno. Dla niespójnego $M$ pierwszą
równość również rozumiemy składowa po składowej
i sumujemy otrzymane miary. W wymiarze $0$ wszystkie
te miary są miarami liczącymi punkty.
\end{twierdzenie}
\begin{proof}
Przypadek wymiaru $0$ wynika z ustalonych konwencji.
W pozostałej części dowodu wymiary są dodatnie.
Zacznijmy od użytego faktu euklidesowego:
przy normalizacji $c_k=\omega_k/2^k$
miara $\mathcal H^k$ na $\R^k$ jest
miarą Lebesgue'a. Wyjaśnijmy rolę stałej.
Nierówność izodiametryczna mówi, że zbiór
o średnicy $D$ ma $k$-wymiarową miarę
zewnętrzną najwyżej $\omega_k(D/2)^k$.
Sumując tę nierówność po dowolnym
pokryciu, otrzymujemy
$\mathcal L^k(A)\leq\mathcal H^k(A)$.
W drugą stronę dla otwartego zbioru
o skończonej objętości wybieramy drobne
kule leżące w nim; twierdzenie Vitalego
o pokryciach wybiera rozłączne kule
pokrywające go z dokładnością do zbioru
miary zero. Ich suma objętości nie
przekracza objętości zbioru. Tu potrzebne jest
osobne uzasadnienie, że reszta zerowej miary Lebesgue'a
ma też zerową miarę Hausdorffa. Dla ustalonych
$\delta,\eta>0$ pokrywamy taki zbiór sześcianami
o bokach $\ell_j<\delta/\sqrt{k}$ i
$\sum_j\ell_j^k<\eta$; pozwala na to definicja miary
zewnętrznej Lebesgue'a, z ewentualnym podziałem sześcianów.
Ich koszt Hausdorffa nie przekracza
$c_k k^{k/2}\eta$. Puszczając $\eta\downarrow0$,
a następnie $\delta\downarrow0$, dostajemy zerowość
bez zakładania dowodzonej równości miar.
Kule z twierdzenia Vitalego wybieramy o średnicach
mniejszych niż $\delta$. Ich koszt Hausdorffa jest
dokładnie sumą objętości, ponieważ $c_k(2r)^k=\omega_k r^k$.
Dodanie pokrycia reszty daje
$\mathcal H^k_\delta(O)\leq\mathcal L^k(O)$ dla otwartego $O$.
Granica po $\delta$ i zewnętrzna regularność miary
Lebesgue'a dają $\mathcal H^k(A)\leq\mathcal L^k(A)$
dla borelowskich $A$ o skończonej mierze.
Dla nieskończonej miary równość wynika już
z przeciwnej nierówności.
Są to klasyczne dwa wyniki analizy miary,
na których opiera się krok euklidesowy.

Przejdźmy do $M$. Ustalmy $0<\varepsilon<1$.
W punkcie $p$ dobierzmy liniowo współrzędne $x$ tak,
by $x(p)=0$ i $G(0)=I_n$. Przy brzegu można to zrobić,
wybierając najpierw ortonormalną bazę przestrzeni stycznej
do brzegu i dopełniając ją normalną do wnętrza;
zmiana liniowa zachowuje model półprzestrzeni.
Na dostatecznie małej domkniętej kuli współrzędnych
$\overline B_{3r}$ (przy brzegu przeciętej z półprzestrzenią)
ciągłość daje
\[
 (1-\varepsilon)|v|^2\leq g_x(v,v)
                         \leq(1+\varepsilon)|v|^2.
\]
Niech $q,q'$ leżą w mniejszej kuli $B_r$.
Odcinek prosty daje oszacowanie górne. Krzywa pozostająca
w $B_{3r}$ ma długość co najmniej
$\sqrt{1-\varepsilon}|x(q)-x(q')|$.
Krzywa wychodząca z $B_{3r}$ musi przed pierwszym wyjściem
przebyć we współrzędnych co najmniej $2r$, podczas gdy
$|x(q)-x(q')|\leq2r$. Dlatego dla \emph{wszystkich}
krzywych, także opuszczających mapę, otrzymujemy
\begin{equation}\label{eq:metryka-lokalnie-blisko-euklidesowej}
 \sqrt{1-\varepsilon}|x(q)-x(q')|
 \leq d_g(q,q')
 \leq\sqrt{1+\varepsilon}|x(q)-x(q')|.
\end{equation}

Dla borelowskiego $A$ zawartego w tej mniejszej mapie
Lemat~\ref{lem:hausdorff-lipschitz} porównuje
$\mathcal H^n_{d_g}(A)$ z $\mathcal L^n(x(A))$
z czynnikami $(1\pm\varepsilon)^{n/2}$.
Te same granice dotyczą $\sqrt{\det G}$, ponieważ
wszystkie wartości własne $G$ leżą w
$[1-\varepsilon,1+\varepsilon]$. Eliminując miarę
Lebesgue'a między oboma porównaniami, dostajemy
\[
 \left(\frac{1-\varepsilon}{1+\varepsilon}\right)^{n/2}\mu_g(A)
 \leq\mathcal H^n_{d_g}(A)
 \leq\left(\frac{1+\varepsilon}{1-\varepsilon}\right)^{n/2}\mu_g(A).
\]
Dla ustalonego $\varepsilon$ wybieramy przeliczalne
pokrycie takimi mapami (przeliczalna baza gwarantuje
istnienie przeliczalnego podpokrycia). Dowolny zbiór
borelowski rozcinamy na rozłączne części: najpierw część
w pierwszej mapie, następnie w drugiej bez pierwszej itd.
Przeliczalna addytywność obu miar daje te same nierówności
globalnie. Dopiero teraz puszczamy $\varepsilon\downarrow0$.
Otrzymujemy równość także dla miary nieskończonej,
gdyż już dowolna dodatnia dolna stała wymusza
nieskończoność drugiej miary.

Pozostaje istotna różnica między odległością otoczenia
a odległością wewnętrzną na $S$. W mapie otoczenia
unormowanej jak wyżej wybierzmy lokalną parametryzację
$F:V\subset\R^k\to S$, a przy brzegu $S$ -- z półprzestrzeni,
tak aby $F(0)=p$ i macierz $J=D(x\circ F)_0$ spełniała
$J^{\mathsf T}J=I_k$. Istnienie takiego wyboru wynika
z iniektywności różniczki i ortonormalizacji bazy $T_pS$.
Zmniejszamy $V$ do wypukłej kuli lub półkuli i dla
$0<\eta<1$ zapewniamy
$\|D(x\circ F)_u-J\|\leq\eta$ na $V$.
Całkując po odcinku od $u$ do $v$, otrzymujemy
\[
 |x(F(u))-x(F(v))-J(u-v)|\leq\eta|u-v|.
\]
W szczególności różnica obrazów ma normę co najmniej
$(1-\eta)|u-v|$. Dobieramy $V$ tak małe, żeby jego
obraz leżał w kuli $B_r$ z
\eqref{eq:metryka-lokalnie-blisko-euklidesowej}.
Każda droga na $S$ jest także drogą na $M$, a obraz
odcinka parametrów jest drogą na $S$. Zatem
\[
 \sqrt{1-\varepsilon}(1-\eta)|u-v|
 \leq d_g(F(u),F(v))
 \leq d_h(F(u),F(v))
 \leq\sqrt{1+\varepsilon}(1+\eta)|u-v|.
\]
Tak samo normy $|D F_u(w)|_g$ są ograniczone przez
te dwie stałe razy $|w|$. Wyznacznik Grama daje więc
dla gęstości $\mu_h$ granice będące $k$-tymi potęgami
tych stałych. Stosujemy poprzednie porównanie miar,
przeliczalne rozcinanie na mapy $S$ oraz
$\varepsilon,\eta\downarrow0$. Obie odległości dają
$\mu_h$, choć globalnie na ogół nie są równe.
Osadzenie zapewnia, że mapy $S$ opisują jego własną
topologię jako podzbioru $M$; nie liczymy wielokrotności
różnych punktów dziedziny immersji.
\end{proof}

Rozróżnijmy dwa użycia twierdzenia:
$\mathcal H^n(M)$ jest objętością całego
$n$-wymiarowego $M$, a
$\mathcal H^{n-1}(\partial M)$ mierzy
pole brzegu. Gdy krzywa sama się przecina,
jej parametr może liczyć fragment
wielokrotnie, podczas gdy miara Hausdorffa
jej \emph{obrazu} liczy każdy punkt raz.
Miara Hausdorffa nie wymaga gładkości;
znaczenie to ma przy zbiorach osobliwych
i granicach rozmaitości.

\section{Dywergencja i twierdzenie Gaussa}
\label{sec:metryka-gauss}

Dywergencja pola wektorowego opisuje lokalne
źródło lub ujście przepływu. Jeżeli przez mały
obszar wypływa więcej niż wpływa, dywergencja
jest dodatnia. Metryka określa zarówno objętość
obszaru, jak i jednostkową normalną potrzebną
do zmierzenia strumienia przez brzeg.

\begin{definicja}[Gradient i dywergencja metryczna]
\index{gradient i dywergencja metryczna@Gradient i dywergencja metryczna}
\label{def:gradient-dywergencja}
Dla $f\in C^\infty(M)$ \emph{gradient} $\nabla_g f$
jest jedynym polem spełniającym
$g(\nabla_g f,Y)=\dd f(Y)$ dla każdego pola $Y$.
Na zorientowanej $(M^n,g)$ \emph{dywergencję}
pola $X$ określa równość
\begin{equation}\label{eq:dywergencja-forma}
 \dd(\iota_X\mathrm{vol}_g)
       =(\operatorname{div}_g X)\,\mathrm{vol}_g.
\end{equation}
Definicja jest niezależna od orientacji:
jej odwrócenie zmienia znak obu form
po obu stronach. Dlatego $\operatorname{div}_gX$
jest zdefiniowana także na rozmaitości
nieorientowalnej: wybieramy orientacje lokalnie, a na
przecięciach ich znaki różnią się o funkcję lokalnie stałą
o wartościach $\pm1$, której różniczka jest zerowa.
\end{definicja}

W każdej przestrzeni stycznej $g_p$ jest
iloczynem skalarnym, więc funkcjonał
$Y\mapsto \dd f_p(Y)$ ma dokładnie jeden
wektor reprezentujący; to uzasadnia
istnienie i jednoznaczność gradientu.
W mapie, przy macierzy odwrotnej
$(g^{ij})=G^{-1}$, otrzymujemy
$(\nabla_gf)^i=\sum_jg^{ij}\partial_jf$.
Współczynniki te są gładkie, bo $G^{-1}$ jest gładka
(wzór z macierzą dopełnień i nieznikającym $\det G$).

\begin{stwierdzenie}[Wzór lokalny i reguła iloczynu]
\label{prop:div-lokalnie}
Jeżeli $X=\sum_i X^i\partial_i$ oraz $G=(g_{ij})$,
to
\begin{equation}\label{eq:dywergencja-lokalna}
 \operatorname{div}_gX
 =\frac{1}{\sqrt{\det G}}\sum_i
          \partial_i\!\left(\sqrt{\det G}\,X^i\right).
\end{equation}
Dla funkcji $f$ zachodzi
$\operatorname{div}_g(fX)
=\dd f(X)+f\operatorname{div}_gX$.
\end{stwierdzenie}
\begin{proof}
W dodatniej mapie połóżmy
$w=\sqrt{\det G}$ i
$\mathrm{vol}_g=w\,\dd x^1\wedge\cdots
\wedge\dd x^n$. Z definicji kontrakcji
\[
 \iota_X\mathrm{vol}_g
 =w\sum_{i=1}^n(-1)^{i-1}X^i\,
    \dd x^1\wedge\cdots\wedge
    \widehat{\dd x^i}\wedge\cdots\wedge\dd x^n.
\]
Po zastosowaniu $\dd$ do $i$-tego składnika
tylko pochodna względem $x^i$ daje
niezerową formę najwyższego stopnia:
każda inna różniczka $\dd x^j$ powtarzałaby
już obecny czynnik. Przeniesienie
$\dd x^i$ na $i$-te miejsce daje kolejny
znak $(-1)^{i-1}$, który kasuje znak
kontrakcji. Otrzymujemy
$\dd(\iota_X\mathrm{vol}_g)
=\sum_i\partial_i(wX^i)\,
  \dd x^1\wedge\cdots\wedge\dd x^n$.
Porównanie z \eqref{eq:dywergencja-forma}
daje \eqref{eq:dywergencja-lokalna}.
Następnie rozwijamy
$\partial_i(wfX^i)=wX^i\partial_i f
  +f\partial_i(wX^i)$ i sumujemy.
\end{proof}

Jednostkowa normalna zewnętrzna istnieje niezależnie
od orientowalności $M$. W mapie przybrzegowej z
$s\geq0$ połóżmy
$\nu=-\nabla_g s/|\nabla_g s|_g$ na $s=0$.
Różniczka $\dd s$ jest niezerowa i znika na wektorach
stycznych do brzegu, więc wzór daje jednostkowy wektor
prostopadły, skierowany na zewnątrz. Inna funkcja
definiująca brzeg ma tam różniczkę $a\,\dd s$, $a>0$;
normalizacja usuwa ten czynnik. Lokalne wzory sklejają
się więc w jedno gładkie pole $\nu$ na $\partial M$.

\begin{twierdzenie}[Twierdzenie Gaussa o dywergencji]
\label{thm:gauss-riemannowski}
Niech $(M^n,g)$ będzie zwartą, zorientowaną
gładką rozmaitością z gładkim brzegiem,
$n\geq1$, a $X$ gładkim polem wektorowym.
Niech $\nu$ będzie na $\partial M$
jednostkową normalną skierowaną na zewnątrz,
$h=i^*g$ metryką brzegu, a
$i:\partial M\hookrightarrow M$ inkluzją.
Wówczas
\begin{equation}\label{eq:gauss-riemannowski}
 \boxed{\displaystyle
 \int_M\operatorname{div}_gX\,\dd\mu_g
 =\int_{\partial M}g(X,\nu)\,\dd\mu_h.}
\end{equation}
Ta sama równość zachodzi dla $M$ niezwartych,
jeśli $X$ ma zwarty nośnik. Obowiązuje również
bez założenia orientowalności, z całkami względem
podanych dodatnich miar.
\end{twierdzenie}
\begin{proof}
Najpierw przyjmijmy orientację. Wstawmy
do Twierdzenia Stokesa~\ref{thm:stokes-formy}
formę $\omega=\iota_X\mathrm{vol}_g$.
Definicja dywergencji daje
\[
 \int_M\operatorname{div}_gX\,\mathrm{vol}_g
 =\int_M\dd\omega
 =\int_{\partial M}i^*\omega.
\]
Dla $n\geq2$ ustalmy punkt brzegu i dodatnią ortonormalną
bazę $(e_1,\ldots,e_{n-1})$ jego przestrzeni
stycznej. Z reguły orientowania brzegu
w Definicji~\ref{def:orientacja-brzegu}
baza $(\nu,e_1,\ldots,e_{n-1})$ jest
dodatnia w $T_pM$; stąd
$\mathrm{vol}_g(\nu,e_1,\ldots,e_{n-1})=1$.
Rozłóżmy $X=g(X,\nu)\nu+X^\top$ na część
normalną i styczną. Forma naprzemienna
zeruje składnik
$\mathrm{vol}_g(X^\top,e_1,\ldots,e_{n-1})$,
bo jest to $n$ wektorów w $(n-1)$-wymiarowej
przestrzeni $T_p\partial M$. Zatem
\[
 (i^*\iota_X\mathrm{vol}_g)
      (e_1,\ldots,e_{n-1})
 =\mathrm{vol}_g(X,e_1,\ldots,e_{n-1})
 =g(X,\nu).
\]
Forma objętościowa brzegu przyjmuje
na $(e_1,\ldots,e_{n-1})$ wartość $1$,
więc $i^*\omega=g(X,\nu)\mathrm{vol}_h$.
Dla $n=1$ używamy konwencji orientacji punktu:
$\mathrm{vol}_h(p)=\mathrm{vol}_g(\nu_p)=\varepsilon_p$.
Ponieważ $X_p=g(X_p,\nu_p)\nu_p$, ta sama tożsamość
$i^*\omega=g(X,\nu)\mathrm{vol}_h$ zachodzi bez wybierania
„dodatniej bazy” przestrzeni zerowej.
To dowodzi \eqref{eq:gauss-riemannowski}
po uwzględnieniu, że dodatnie formy
odpowiadają miarom $\mu_g$ i $\mu_h$.

Dla nieorientowalnego $M$ bierzemy podział
jedności $(\rho_j)$ podporządkowany mapom
z lokalnie wybraną orientacją i stosujemy
powyższy dowód do pól $\rho_jX$ w tych mapach.
Ich nośniki nie dotykają sztucznych brzegów map.
Suma lewych stron jest całką z
$\sum_j\operatorname{div}_g(\rho_jX)$,
a reguła iloczynu daje
\[
 \sum_j\operatorname{div}_g(\rho_jX)
 =\left(\sum_j\rho_j\right)\operatorname{div}_gX
  +\dd\!\left(\sum_j\rho_j\right)(X)
 =\operatorname{div}_gX.
\]
Na rzeczywistym brzegu
$\sum_jg(\rho_jX,\nu)=g(X,\nu)$.
Lokalne wybory orientacji nie wpływają
na całki z dodatnich gęstości.
Ten sam argument z lokalnie skończonym
podziałem działa dla pola o zwartym nośniku
na niezwartym $M$; po pomnożeniu przez $X$
pozostaje tylko skończenie wiele składników.
\end{proof}

\begin{wniosek}[Całkowanie przez części]
\label{cor:riemann-calkowanie-przez-czesci}
Przy założeniach Twierdzenia
\ref{thm:gauss-riemannowski}, dla gładkiej
funkcji $f$ i pola $X$,
\[
 \int_M g(\nabla_gf,X)\,\dd\mu_g
 =-\int_M f\operatorname{div}_gX\,\dd\mu_g
   +\int_{\partial M}f\,g(X,\nu)\,\dd\mu_h.
\]
Jeśli $\partial M=\varnothing$, to
$\int_M\operatorname{div}_gX\,\dd\mu_g=0$.
\end{wniosek}
\begin{proof}
Z definicji gradientu $\dd f(X)=g(\nabla_gf,X)$.
Zastosujmy \eqref{eq:gauss-riemannowski}
do $fX$ i podstawmy regułę iloczynu
ze Stwierdzenia~\ref{prop:div-lokalnie}.
Przenieśmy składnik z $f\operatorname{div}X$
na prawą stronę. Dla drugiej tezy podstawmy
$f=1$ i pusty brzeg.
\end{proof}

\begin{przyklad}[Kula euklidesowa i strumień]
Dla $M=\overline B_R(0)\subset\R^3$
i $X(x)=x$ wzór \eqref{eq:dywergencja-lokalna}
daje $\operatorname{div}X=3$.
Na sferze $\nu=x/R$, więc $X\cdot\nu=R$.
Obie strony twierdzenia Gaussa są równe
$3(4\pi R^3/3)=R(4\pi R^2)=4\pi R^3$.
Odzyskujemy klasyczne Twierdzenie
Gaussa--Ostrogradskiego
z Wniosku~\ref{cor:stokes-gauss}.

Dla wykresu $z=f(x,y)$ zorientowanego ku górze
wektor normalny i element pola wynoszą
\[
 \nu=\frac{(-f_x,-f_y,1)}
                 {\sqrt{1+f_x^2+f_y^2}},
 \qquad
 \dd A=\sqrt{1+f_x^2+f_y^2}\,\dd x\,\dd y.
\]
Jeżeli $X=(0,0,1)$, to
$g(X,\nu)\dd A=\dd x\,\dd y$.
Strumień pola pionowego przez wykres nad
dyskiem promienia $a$ wynosi więc $\pi a^2$,
mimo że samo pole zakrzywionej powierzchni
może być od $\pi a^2$ większe.
\end{przyklad}

\par\smallskip
{\footnotesize\noindent\emph{Dalsza lektura:}
Trzy modele hiperboliczne i ich izometrie omawiają
\href{https://math.uchicago.edu/~dannyc/courses/kleinian_groups_2015/kleinian_groups.pdf}
{notatki Danny'ego Calegariego o geometrii hiperbolicznej}.
Postacie lorentzowskich metryk
\href{https://sites.science.oregonstate.edu/physics/coursewikis/GGR/_export/xhtml/book/ggr/schwarz.html}
{Schwarzschilda} i
\href{https://sites.science.oregonstate.edu/coursewikis/GGR/book/content/kerr}
{Kerra} opisuje kurs \emph{Geometry of General Relativity}
Uniwersytetu Stanowego Oregonu.
Równość miary Hausdorffa i objętości
riemannowskiej wraz z formułą pola rozwija
\href{https://doi.org/10.1017/CBO9780511616822.005}
{Isaac Chavel, \emph{Riemannian Geometry:
A Modern Introduction}, część III}.
Fakty o miarze Hausdorffa w przestrzeni euklidesowej,
w tym nierówność izodiametryczną i pokrycia Vitalego,
omawia dodatek A do
\href{https://www.math.cmu.edu/~gautam/sj/teaching/2022-23/720-measure/pdfs/measure.pdf}
{\emph{Lecture Notes on Measure Theory} Gautama Iyera}.\par}
